% Part: proof-theory
% Chapter: sequent-calculus
% Section: rules-proofs

\documentclass[../../../include/open-logic-section]{subfiles}

\begin{document}

\olfileid{pt}{seq}{rp}

\olsection{Aturan lan \usetoken{P}{proof}}

Kita miwiti kanthi menehi sawetara definisi lan netepake tetembungan.

Contone, gatekna loro aturan ing \Log{G3c} kang ngidini !!{proving}
sekuen kang ngemot $!A \land !B$ saka sekuen kang ngemot $!A$
lan~$!B$:
\[
\Axiom$ !A, !B, \Gamma \fCenter \Delta$
\RightLabel{\LeftR{\land}}
\UnaryInf$ !A \land !B, \Gamma \fCenter \Delta$
\DisplayProof
\qquad
\Axiom$\Gamma \fCenter \Delta, !A$
\Axiom$ \Gamma \fCenter \Delta, !B$
\RightLabel{\RightR{\land}}
\BinaryInf$ \Gamma \fCenter \Delta, !A \land !B $
\DisplayProof
\]
Sekuen ing ndhuwur garis diarani \emph{premis}, lan sekuen ing
ngisor garis diarani \emph{kesimpulan}. !!{formula}s ing
$\Gamma$ lan $\Delta$ ing saben aturan diarani \emph{konteks} utawa
\emph{!!{formula}s sisih}. !!{formula} ing kesimpulan kang
dudu bagean saka konteks diarani !!{formula} \emph{utama},
lan !!{formula}s ing premis kang dudu bagean saka konteks diarani
!!{formula}s \emph{aktif} (utawa \emph{pambiyantu}).

Aturan kuantor rada luwih rumit. Contone, aturan ing
\Log{G1c} kanggo $\lforall$ yaiku:
\[\Axiom$ !A(t), \Gamma \fCenter \Delta$
\RightLabel{\LeftR{\lforall}}
\UnaryInf$ \lforall[x][!A(x)],\Gamma \fCenter \Delta$
\DisplayProof
\qquad
\Axiom$ \Gamma \fCenter \Delta, !A(a) $
\RightLabel{\RightR{\lforall}}
\UnaryInf$ \Gamma \fCenter \Delta, \lforall[x][!A(x)]$
\DisplayProof
\]
Ing aturan \LeftR{\lforall}, !!{formula} aktif yaiku~$!A(t)$,
kanthi $t$~minangka term katutup, yaiku term tanpa variabel.
Inferensi kaya mangkono bener---tegese cocog karo aturan skematis
\LeftR{\lforall}---yen ana !!{formula}~$!A$,
variabel~$x$, lan term katutup~$t$ kang ndadekake
!!{formula} utama ing kesimpulan minangka $\lforall[x][!A]$ lan
!!{formula} aktif ing premis minangka~$\Subst{!A}{t}{x}$.
Aturan \RightR{\lforall} lumaku kanthi cara kang padha, nanging
!!{formula} aktif ing premis ora nggunakake sembarang term~$t$;
formula kasebut kudu awujud~$\Subst{!A}{a}{x}$,
kanthi $a$ minangka !!a{constant}. Konstanta iki
diarani \emph{eigenvariabel} saka inferensi \RightR{\lforall}.
Inferensi kasebut mung bener yen sekuen ing ngisor ora ngemot~$a$,
yaiku yen $a$ ora ana ing $\Gamma$, $\Delta$, utawa~$!A$.
Iki diarani \emph{syarat eigenvariabel}.

Sistem sekuen kayata \Log{G1c} lan \Log{G3c} dumadi saka kumpulan
aturan kang ditemtokake kanthi skematis kaya mangkono, bebarengan karo
sawetara sekuen kang dianggep aksioma lan uga ditemtokake kanthi skematis.
!!^{proof}s ing sistem kasebut minangka wit sekuen winates kang
dibangun kanthi induktif saka aksioma nganggo aturan inferensi.
Saben godhong minangka aksioma, lan saben sekuen liyane asale saka
sekuen kang ana langsung ing ndhuwure kanthi salah siji aturan inferensi
ing sistem kasebut.

\begin{defn}\ollabel{defn:proof}
\emph{!!^{proof}s} ing sistem sekuen minangka wit sekuen winates
kang dibangun kanthi induktif kaya mangkene:
\begin{enumerate}
  \item\ollabel{defn:prf:axiom} Saben aksioma minangka !!a{proof}.
  \item\ollabel{defn:prf:one-prem} Yen $\pi_1$ minangka !!a{proof}
  kang dipungkasi dening sekuen~$S_1$, lan $S$ minangka sekuen
  kang asale saka~$S_1$ nganggo aturan~$R$ kanthi siji premis, mula
  \[
  \AxiomC{}
  \RightLabel{$\pi_1$}
  \DeduceC{$S_1$}
  \RightLabel{$R$}
  \UnaryInfC{$S$}
  \DisplayProof
  \]
  minangka !!a{proof}.
  \item\ollabel{defn:prf:two-prem} Yen $\pi_1$ lan $\pi_2$ minangka
  !!{proof}s kang dipungkasi dening sekuen~$S_1$ lan~$S_2$, lan
  $S$ minangka sekuen kang asale saka~$S_1$ lan $S_2$ nganggo
  aturan~$R$ kanthi loro premis, mula
  \[
  \AxiomC{}
  \RightLabel{$\pi_1$}
  \DeduceC{$S_1$}
  \AxiomC{}
  \RightLabel{$\pi_2$}
  \DeduceC{$S_2$}
  \RightLabel{$R$}
  \BinaryInfC{$S$}
  \DisplayProof
  \]
  minangka !!a{proof}.
\end{enumerate}
\end{defn}

\Olref{defn:proof} nemtokake !!{proof}s minangka wit sekuen.
Kita nganggep wit kasebut tuwuh ``munggah.'' Simpul paling ndhuwur
(godhong) ing wit iki minangka aksioma lan diarani \emph{sekuen wiwitan}.
Sekuen paling ngisor diarani \emph{sekuen pungkasan}. Yen $\pi$
minangka !!a{proof} kanthi sekuen pungkasan~$S$, kita uga nulis
$\pi \Proves S$. Yen ana !!a{proof} kanggo sekuen~$S$, kita nulis
$\Proves S$. Sistem kang ngemot bukti kasebut bisa dituduhake ing
sisih kiwa tandha $\Proves$, contone,
$\Log{G1c} \Proves !A, !B \Sequent !A \land !B$.
Saben sekuen ing !!a{proof}, kajaba sekuen wiwitan, minangka
\emph{kesimpulan} saka inferensi nganggo sawijining aturan~$R$.
Siji utawa loro sekuen kang ana langsung ing ndhuwur sawijining sekuen
ing !!a{proof} (yaiku simpul wong tuwane) diarani \emph{premis}
saka inferensi kasebut.

!!{proof}s kang digunakake ing pambangunan induktif sawijining
!!a{proof} diarani \emph{!!{proof}s bagean} saka bukti kasebut.
!!{proof}s bagean kang dipungkasi ing premis sawijining inferensi
diarani \emph{!!{proof}s bagean langsung} saka inferensi kasebut.

Asring perlu ngukur kerumitan !!{proof}s nganggo wilangan asli.
Kita bisa mung ngetung cacah sekuen ing !!a{proof}, nanging ukuran
kang luwih migunani asring yaiku !!{height} saka !!a{proof}:
yaiku cacah langkah inferensi paling gedhe ing sawijining jalur
saka sekuen pungkasan menyang sekuen wiwitan. Kita uga bisa menehi
definisi induktif.

\begin{defn}
\emph{!!{height}}~$\pheight{\pi}$ saka !!a{proof}~$\pi$
ditemtokake kanthi induktif kaya mangkene.
\begin{enumerate}
  \item Yen $\pi$ mung dumadi saka siji aksioma, $\pheight{\pi} = 0$.
  \item Yen $\pi$ dipungkasi dening inferensi kanthi siji premis lan
  !!{proof} bagean langsung~$\pi_1$, mula
  $\pheight{\pi} = \pheight{\pi_1} + 1$.
  \item Yen $\pi$ dipungkasi dening inferensi kanthi loro premis lan
  !!{proof}s bagean langsung~$\pi_1$ lan~$\pi_2$, mula
  $\pheight{\pi} = \max(\pheight{\pi_1}, \pheight{\pi_2}) + 1$.
\end{enumerate}
Yen sekuen~$S$ nduweni !!a{proof} kanthi dhuwur $\le n$,
kita nulis $\Proves[n] S$.
\end{defn}

\begin{ex}
Ing \Log{G3c}, saben sekuen awujud $!C, \Gamma \Sequent \Delta, !C$
minangka aksioma, kanthi $!C$ kudu atomik. Upamane $!D$ lan $!E$
minangka ukara atomik. Mula $!D, !E \Sequent !D$ lan
$!D, !E \Sequent !E$ minangka instans saka aksioma
$!C, \Gamma \Sequent \Delta, !C$. Contone, yen kita milih
$!C \ident !E$, $\Gamma = \{!D\}$ lan $\Delta = \emptyset$ ing
$!C, \Gamma \Sequent \Delta, !C$, kita oleh sekuen
$!D, !E \Sequent !E$ kanthi pas. (Elinga yen kita nggunakake
multi\emph{set}, mula $!D, !E \Sequent !E$ lan
$!E, !D \Sequent !E$ minangka sekuen kang padha.)

Amarga $!D, !E \Sequent !D$ minangka instans saka aksioma
ing~\Log{G3c}, sekuen kasebut dhewe dianggep !!a{proof}
ing~\Log{G3c}. Semono uga $!D, !E \Sequent !E$, miturut
\olref{defn:proof}\olref{defn:prf:axiom}. Ayo diarani
$\pi_1$ lan~$\pi_2$. Miturut \olref{defn:prf:two-prem}, kita bisa
njupuk loro !!{proof}s kasebut lan mbangun bukti anyar ($\pi_3$)
nganggo aturan~$\RightR{\land}$ kanthi loro premis:
\[
\Axiom$!D, !E \fCenter !D$
\Axiom$!D, !E \fCenter !E$
\RightLabel{\RightR{\land}}
\BinaryInf$!D, !E  \fCenter !D \land !E $
\DisplayProof
\]
Saka $\pi_3$, sabanjure kita oleh !!{proof}~$\pi_4$
\[
\Axiom$!D, !E \fCenter !D$
\Axiom$!D, !E \fCenter !E$
\RightLabel{\RightR{\land}}
\insertBetweenHyps{\hspace{5ex}}
\BinaryInf$!D, !E  \fCenter !D \land !E $
\LeftSubproofLabel{$\pi_3$}
\RightLabel{\LeftR{\land}}
\UnaryInf$!D \land !E  \fCenter !D \land !E$
\RightSubproofLabel{$\pi_4$}
\DisplayProof
\]
nganggo aturan~$\LeftR{\land}$ kanthi siji premis (nggunakake
\olref{defn:proof}\olref{defn:prf:one-prem}).
!!{proof} bagean langsung saka inferensi pungkasan ing~$\pi_4$
yaiku~$\pi_3$. !!{proof}s bagean langsung saka inferensi
$\RightR{\land}$ yaiku $\pi_1$ lan~$\pi_2$.
!!{proof}s bagean saka $\pi_4$ yaiku kabeh $\pi_1$, $\pi_2$,
$\pi_3$, lan $\pi_4$ dhewe. Kita nduweni $\pheight{\pi_1} =
\pheight{\pi_2} = 0$, $\pheight{\pi_3} = 1$ lan $\pheight{\pi_4} = 2$.
Kita wis nuduhake yen $\Log{G3c} \Proves[4] !D \land !E  \fCenter !D
\land !E$ kanggo sembarang $!D$ lan~$!E$ atomik.
\end{ex}

\end{document}
